\section{GPU Acceleration and CUDA Interpolation}

\subsection{Why the Grid Needs a GPU}

The propagation grid has of order $10^3$ radial points, $10^3$ to $4\times10^3$ time points, and $10^4$ propagation steps, so the field envelope alone is an array of several million complex numbers that is updated at every step. Every step also needs two matrix multiplications with the QDHT matrix $Y$, several FFTs along the time axis, and a full evaluation of the nonlinear right hand side at four Runge-Kutta stages. On a CPU this would take hours to days for a single run. The code offloads all of these operations to the GPU through \texttt{cupy}, which is a drop in replacement for \texttt{numpy} that runs the same array operations on CUDA hardware, so the matrix products, the FFTs and the elementwise nonlinear evaluations all execute in parallel across thousands of GPU cores instead of sequentially. In practice the $4\ \mu$J run used in this report takes about four minutes for $9000$ steps.

\subsection{The Photoionization Rate Table}

The Keldysh photoionization rate $W_{PI}$ derived in Section~3 involves an infinite sum of Dawson functions and complete elliptic integrals. This is far too expensive to evaluate at every grid point and every time step during propagation, since it would dominate the runtime. Instead the code evaluates $W_{PI}$ once, before propagation starts, on a dense logarithmically spaced grid of intensities spanning the full range the simulation can reach, and fits a monotone cubic Hermite interpolant through these points. The result is a lookup table of $4096$ points in $\log_{10} I$, stored on the GPU. There are two such tables, one for the valence to conduction band gap $U_i$ and one for the shallower self-trapped exciton gap $U_s$, which is what makes the re-ionization term of Section~11.2 cheap enough to evaluate at every time sample.

During propagation, whenever the code needs $W_{PI}$ at some intensity, it no longer evaluates the analytical Keldysh formula. It reads the table and interpolates linearly between the two nearest samples in log intensity space. Because the table is sampled logarithmically and the Keldysh rate varies smoothly on a logarithmic intensity scale, this turns a computation involving dozens of transcendental function evaluations into two table reads and one multiplication, at the cost of a small and controlled interpolation error.

\subsection{The CUDA Kernel for the Rate Equations}

The interpolation above is not only used from the array level. It is also compiled directly into a custom CUDA kernel, written as a raw CUDA C string in \texttt{sim/kernels\dotb py} and launched through \texttt{cupy\dotb RawKernel}. This is a deliberate departure from the array based style used everywhere else in the code, and it exists because the rate equations of Section~11.2 have to be integrated sequentially in time at every radius. One time sample depends on the previous one, so this step cannot be vectorized as a single array operation the way the FFTs and the QDHT product can.

The kernel assigns one CUDA thread to each radial point. Each thread then walks the entire time axis in a tight loop, from the first sample to the last, evaluating the interpolated rate at every step and advancing $N$ and $N_{STE}$ with the exact exponential update described in Section~12.3. Because thousands of radial points are independent of one another, thousands of these sequential time loops run at the same time on the GPU, one per thread. This is the usual way to make an inherently sequential recurrence fast, which is to parallelize across the dimension that has no sequential dependency, here the radius, instead of trying to parallelize the one that does, here the time.

Two details of the kernel are worth pointing out because they protect the physics rather than the performance. The first sample of every radial line is reset to zero before the loop starts, so each propagation step begins from an unexcited medium at the front of the temporal window instead of inheriting a leftover value. The second is the rescaling applied when $N+N_{STE}$ exceeds $N_{at}$, which keeps the two populations from together claiming more sites than the material has. Both are cheap guards inside the inner loop, and both matter near the clamping intensity where the depletion factor is doing real work.

\subsection{Summary of the GPU Division of Labour}

Three different accelerations coexist in the code, each solving a different bottleneck. The QDHT matrix product and the FFTs are handled by the built in \texttt{cupy} linear algebra and FFT routines, which are already GPU native and need no custom code. The photoionization rate is precomputed once on the CPU with \texttt{scipy}, then transferred to the GPU as a lookup table, which trades an expensive repeated computation for a cheap repeated memory read. The population rate equations, which are sequential in time and cannot be written as a single array operation, are handled by the hand written CUDA kernel that exploits parallelism across the radial coordinate instead. Together these three choices are what make it possible to run a full simulation in two transverse dimensions plus time, over tens of thousands of propagation steps, in minutes rather than days.
